% Options for packages loaded elsewhere
\PassOptionsToPackage{unicode}{hyperref}
\PassOptionsToPackage{hyphens}{url}
%
\documentclass[
]{article}
\usepackage{amsmath,amssymb}
\usepackage{lmodern}
\usepackage{iftex}
\ifPDFTeX
  \usepackage[T1]{fontenc}
  \usepackage[utf8]{inputenc}
  \usepackage{textcomp} % provide euro and other symbols
\else % if luatex or xetex
  \usepackage{unicode-math}
  \defaultfontfeatures{Scale=MatchLowercase}
  \defaultfontfeatures[\rmfamily]{Ligatures=TeX,Scale=1}
\fi
% Use upquote if available, for straight quotes in verbatim environments
\IfFileExists{upquote.sty}{\usepackage{upquote}}{}
\IfFileExists{microtype.sty}{% use microtype if available
  \usepackage[]{microtype}
  \UseMicrotypeSet[protrusion]{basicmath} % disable protrusion for tt fonts
}{}
\makeatletter
\@ifundefined{KOMAClassName}{% if non-KOMA class
  \IfFileExists{parskip.sty}{%
    \usepackage{parskip}
  }{% else
    \setlength{\parindent}{0pt}
    \setlength{\parskip}{6pt plus 2pt minus 1pt}}
}{% if KOMA class
  \KOMAoptions{parskip=half}}
\makeatother
\usepackage{xcolor}
\IfFileExists{xurl.sty}{\usepackage{xurl}}{} % add URL line breaks if available
\IfFileExists{bookmark.sty}{\usepackage{bookmark}}{\usepackage{hyperref}}
\hypersetup{
  pdftitle={Session-0.-Introduccion a R y R Studio.},
  pdfauthor={Martin Murillo},
  hidelinks,
  pdfcreator={LaTeX via pandoc}}
\urlstyle{same} % disable monospaced font for URLs
\usepackage[margin=1in]{geometry}
\usepackage{color}
\usepackage{fancyvrb}
\newcommand{\VerbBar}{|}
\newcommand{\VERB}{\Verb[commandchars=\\\{\}]}
\DefineVerbatimEnvironment{Highlighting}{Verbatim}{commandchars=\\\{\}}
% Add ',fontsize=\small' for more characters per line
\usepackage{framed}
\definecolor{shadecolor}{RGB}{248,248,248}
\newenvironment{Shaded}{\begin{snugshade}}{\end{snugshade}}
\newcommand{\AlertTok}[1]{\textcolor[rgb]{0.94,0.16,0.16}{#1}}
\newcommand{\AnnotationTok}[1]{\textcolor[rgb]{0.56,0.35,0.01}{\textbf{\textit{#1}}}}
\newcommand{\AttributeTok}[1]{\textcolor[rgb]{0.77,0.63,0.00}{#1}}
\newcommand{\BaseNTok}[1]{\textcolor[rgb]{0.00,0.00,0.81}{#1}}
\newcommand{\BuiltInTok}[1]{#1}
\newcommand{\CharTok}[1]{\textcolor[rgb]{0.31,0.60,0.02}{#1}}
\newcommand{\CommentTok}[1]{\textcolor[rgb]{0.56,0.35,0.01}{\textit{#1}}}
\newcommand{\CommentVarTok}[1]{\textcolor[rgb]{0.56,0.35,0.01}{\textbf{\textit{#1}}}}
\newcommand{\ConstantTok}[1]{\textcolor[rgb]{0.00,0.00,0.00}{#1}}
\newcommand{\ControlFlowTok}[1]{\textcolor[rgb]{0.13,0.29,0.53}{\textbf{#1}}}
\newcommand{\DataTypeTok}[1]{\textcolor[rgb]{0.13,0.29,0.53}{#1}}
\newcommand{\DecValTok}[1]{\textcolor[rgb]{0.00,0.00,0.81}{#1}}
\newcommand{\DocumentationTok}[1]{\textcolor[rgb]{0.56,0.35,0.01}{\textbf{\textit{#1}}}}
\newcommand{\ErrorTok}[1]{\textcolor[rgb]{0.64,0.00,0.00}{\textbf{#1}}}
\newcommand{\ExtensionTok}[1]{#1}
\newcommand{\FloatTok}[1]{\textcolor[rgb]{0.00,0.00,0.81}{#1}}
\newcommand{\FunctionTok}[1]{\textcolor[rgb]{0.00,0.00,0.00}{#1}}
\newcommand{\ImportTok}[1]{#1}
\newcommand{\InformationTok}[1]{\textcolor[rgb]{0.56,0.35,0.01}{\textbf{\textit{#1}}}}
\newcommand{\KeywordTok}[1]{\textcolor[rgb]{0.13,0.29,0.53}{\textbf{#1}}}
\newcommand{\NormalTok}[1]{#1}
\newcommand{\OperatorTok}[1]{\textcolor[rgb]{0.81,0.36,0.00}{\textbf{#1}}}
\newcommand{\OtherTok}[1]{\textcolor[rgb]{0.56,0.35,0.01}{#1}}
\newcommand{\PreprocessorTok}[1]{\textcolor[rgb]{0.56,0.35,0.01}{\textit{#1}}}
\newcommand{\RegionMarkerTok}[1]{#1}
\newcommand{\SpecialCharTok}[1]{\textcolor[rgb]{0.00,0.00,0.00}{#1}}
\newcommand{\SpecialStringTok}[1]{\textcolor[rgb]{0.31,0.60,0.02}{#1}}
\newcommand{\StringTok}[1]{\textcolor[rgb]{0.31,0.60,0.02}{#1}}
\newcommand{\VariableTok}[1]{\textcolor[rgb]{0.00,0.00,0.00}{#1}}
\newcommand{\VerbatimStringTok}[1]{\textcolor[rgb]{0.31,0.60,0.02}{#1}}
\newcommand{\WarningTok}[1]{\textcolor[rgb]{0.56,0.35,0.01}{\textbf{\textit{#1}}}}
\usepackage{graphicx}
\makeatletter
\def\maxwidth{\ifdim\Gin@nat@width>\linewidth\linewidth\else\Gin@nat@width\fi}
\def\maxheight{\ifdim\Gin@nat@height>\textheight\textheight\else\Gin@nat@height\fi}
\makeatother
% Scale images if necessary, so that they will not overflow the page
% margins by default, and it is still possible to overwrite the defaults
% using explicit options in \includegraphics[width, height, ...]{}
\setkeys{Gin}{width=\maxwidth,height=\maxheight,keepaspectratio}
% Set default figure placement to htbp
\makeatletter
\def\fps@figure{htbp}
\makeatother
\setlength{\emergencystretch}{3em} % prevent overfull lines
\providecommand{\tightlist}{%
  \setlength{\itemsep}{0pt}\setlength{\parskip}{0pt}}
\setcounter{secnumdepth}{-\maxdimen} % remove section numbering
\ifLuaTeX
  \usepackage{selnolig}  % disable illegal ligatures
\fi

\title{Session-0.-Introduccion a R y R Studio.}
\author{Martin Murillo}
\date{2026-05-16}

\begin{document}
\maketitle

\hypertarget{introduccion}{%
\section{INTRODUCCION}\label{introduccion}}

\hypertarget{el-presente-cuxf3digo-en-lenguaje-r-tiene-como-objetivo-analizar-paruxe1metros-ambientales-y-fisicoquuxedmicos-registrados-en-diferentes-transectos-de-un-ruxedo-durante-tres-periodos-de-muestreo-mayo-septiembre-y-octubre.-para-ello-se-emplean-distintas-libreruxedas-de-r-enfocadas-en-manejo-de-datos-anuxe1lisis-estaduxedstico-y-visualizaciuxf3n-gruxe1fica.}{%
\subsection{El presente código en lenguaje R tiene como objetivo
analizar parámetros ambientales y fisicoquímicos registrados en
diferentes transectos de un río durante tres periodos de muestreo (mayo,
septiembre y octubre). Para ello, se emplean distintas librerías de R
enfocadas en manejo de datos, análisis estadístico y visualización
gráfica.}\label{el-presente-cuxf3digo-en-lenguaje-r-tiene-como-objetivo-analizar-paruxe1metros-ambientales-y-fisicoquuxedmicos-registrados-en-diferentes-transectos-de-un-ruxedo-durante-tres-periodos-de-muestreo-mayo-septiembre-y-octubre.-para-ello-se-emplean-distintas-libreruxedas-de-r-enfocadas-en-manejo-de-datos-anuxe1lisis-estaduxedstico-y-visualizaciuxf3n-gruxe1fica.}}

\hypertarget{el-script-inicia-con-la-carga-y-organizaciuxf3n-de-bases-de-datos-en-formato-.csv-posteriormente-se-extraen-variables-ambientales-como-temperatura-conductividad-oxuxedgeno-disuelto-ph-humedad-velocidad-de-corriente-y-caudal.-a-partir-de-estos-datos-se-calculan-medidas-de-tendencia-central-especuxedficamente-promedios-mensuales-con-el-fin-de-comparar-las-condiciones-ambientales-entre-los-diferentes-periodos-de-estudio.}{%
\subsection{El script inicia con la carga y organización de bases de
datos en formato .csv, posteriormente se extraen variables ambientales
como temperatura, conductividad, oxígeno disuelto, pH, humedad,
velocidad de corriente y caudal. A partir de estos datos, se calculan
medidas de tendencia central, específicamente promedios mensuales, con
el fin de comparar las condiciones ambientales entre los diferentes
periodos de
estudio.}\label{el-script-inicia-con-la-carga-y-organizaciuxf3n-de-bases-de-datos-en-formato-.csv-posteriormente-se-extraen-variables-ambientales-como-temperatura-conductividad-oxuxedgeno-disuelto-ph-humedad-velocidad-de-corriente-y-caudal.-a-partir-de-estos-datos-se-calculan-medidas-de-tendencia-central-especuxedficamente-promedios-mensuales-con-el-fin-de-comparar-las-condiciones-ambientales-entre-los-diferentes-periodos-de-estudio.}}

\hypertarget{finalmente-el-cuxf3digo-integra-los-resultados-en-tablas-resumen-y-desarrolla-un-anuxe1lisis-de-agrupamiento-jeruxe1rquico-cluster-analysis-utilizando-distancia-euclidiana-permitiendo-evaluar-la-similitud-entre-meses-seguxfan-los-paruxe1metros-ambientales-registrados.-los-resultados-obtenidos-se-exportan-en-formatos-.csv-y-gruxe1ficos-.png-para-su-posterior-interpretaciuxf3n-y-presentaciuxf3n.}{%
\subsection{Finalmente, el código integra los resultados en tablas
resumen y desarrolla un análisis de agrupamiento jerárquico (cluster
analysis) utilizando distancia euclidiana, permitiendo evaluar la
similitud entre meses según los parámetros ambientales registrados. Los
resultados obtenidos se exportan en formatos .csv y gráficos .png para
su posterior interpretación y
presentación.}\label{finalmente-el-cuxf3digo-integra-los-resultados-en-tablas-resumen-y-desarrolla-un-anuxe1lisis-de-agrupamiento-jeruxe1rquico-cluster-analysis-utilizando-distancia-euclidiana-permitiendo-evaluar-la-similitud-entre-meses-seguxfan-los-paruxe1metros-ambientales-registrados.-los-resultados-obtenidos-se-exportan-en-formatos-.csv-y-gruxe1ficos-.png-para-su-posterior-interpretaciuxf3n-y-presentaciuxf3n.}}

\hypertarget{establecer-parametros-de-paquete-knit-y-carga-de-paquetes-en-r}{%
\section{Establecer parametros de paquete knit y carga de paquetes en
R}\label{establecer-parametros-de-paquete-knit-y-carga-de-paquetes-en-r}}

\begin{Shaded}
\begin{Highlighting}[]
\NormalTok{knitr}\SpecialCharTok{::}\NormalTok{opts\_chunk}\SpecialCharTok{$}\FunctionTok{set}\NormalTok{(}\AttributeTok{echo =} \ConstantTok{TRUE}\NormalTok{, }\AttributeTok{cahe=}\ConstantTok{TRUE}\NormalTok{, }\AttributeTok{message=}\ConstantTok{FALSE}\NormalTok{, }\AttributeTok{warning=}\ConstantTok{FALSE}\NormalTok{, }\AttributeTok{fig.width=}\DecValTok{12}\NormalTok{, }\AttributeTok{fig.height=}\FloatTok{4.5}\NormalTok{, }\AttributeTok{error=}\ConstantTok{TRUE}\NormalTok{,}\AttributeTok{cache =} \ConstantTok{TRUE}\NormalTok{)}
\CommentTok{\# Establecer formatos de agrupacion de codigos de R}

\CommentTok{\# Cargar paquetes de R }
\NormalTok{xfun}\SpecialCharTok{::}\FunctionTok{pkg\_attach}\NormalTok{(}\StringTok{"readr"}\NormalTok{, }\StringTok{"here"}\NormalTok{, }\StringTok{"dplyr"}\NormalTok{, }\StringTok{"tidyr"}\NormalTok{, }\StringTok{"tibble"}\NormalTok{, }\StringTok{"plotrix"}\NormalTok{,}\StringTok{"ggplot2"}\NormalTok{, }\StringTok{"Hmisc"}\NormalTok{, }\StringTok{"readxl"}\NormalTok{,}\StringTok{"unmarked"}\NormalTok{,}\StringTok{"AHMbook"}\NormalTok{, }\StringTok{"AICcmodavg"}\NormalTok{, }\StringTok{"tidyverse"}\NormalTok{, }\StringTok{"vegan"}\NormalTok{, }\StringTok{"DHARMa"}\NormalTok{, }\StringTok{"corrplot"}\NormalTok{, }\StringTok{"kimisc"}\NormalTok{, }\StringTok{"car"}\NormalTok{, }\StringTok{"abdiv"}\NormalTok{, }\StringTok{"ggfortify"}\NormalTok{, }\StringTok{"ggpubr"}\NormalTok{, }\StringTok{"qwraps2"}\NormalTok{, }\StringTok{"car"}\NormalTok{)}
\end{Highlighting}
\end{Shaded}

\begin{verbatim}
## here() starts at C:/R/Asesorias/Curso_E_P_C_2026/Materiales_Iniciales/Repositorio_GitHub/Session-0.-Introducci-n-a-R-y-R-Studio
\end{verbatim}

\begin{verbatim}
## 
## Attaching package: 'dplyr'
\end{verbatim}

\begin{verbatim}
## The following objects are masked from 'package:stats':
## 
##     filter, lag
\end{verbatim}

\begin{verbatim}
## The following objects are masked from 'package:base':
## 
##     intersect, setdiff, setequal, union
\end{verbatim}

\begin{verbatim}
## 
## Attaching package: 'Hmisc'
\end{verbatim}

\begin{verbatim}
## The following objects are masked from 'package:dplyr':
## 
##     src, summarize
\end{verbatim}

\begin{verbatim}
## The following objects are masked from 'package:base':
## 
##     format.pval, units
\end{verbatim}

\begin{verbatim}
## -- Attaching core tidyverse packages ------------------------ tidyverse 2.0.0 --
## v forcats   1.0.0     v purrr     1.0.4
## v lubridate 1.9.4     v stringr   1.5.2
## -- Conflicts ------------------------------------------ tidyverse_conflicts() --
## x dplyr::filter()    masks stats::filter()
## x dplyr::lag()       masks stats::lag()
## x Hmisc::src()       masks dplyr::src()
## x Hmisc::summarize() masks dplyr::summarize()
## i Use the conflicted package (<http://conflicted.r-lib.org/>) to force all conflicts to become errors
## Loading required package: permute
## 
## This is DHARMa 0.4.7. For overview type '?DHARMa'. For recent changes, type news(package = 'DHARMa')
## 
## corrplot 0.95 loaded
## 
## Loading required package: carData
## 
## 
## Attaching package: 'car'
## 
## 
## The following object is masked from 'package:purrr':
## 
##     some
## 
## 
## The following object is masked from 'package:unmarked':
## 
##     vif
## 
## 
## The following object is masked from 'package:dplyr':
## 
##     recode
## 
## 
## 
## Attaching package: 'abdiv'
## 
## 
## The following object is masked from 'package:unmarked':
## 
##     richness
## 
## 
## 
## Attaching package: 'qwraps2'
## 
## 
## The following objects are masked from 'package:ggpubr':
## 
##     mean_ci, mean_sd, median_iqr
## 
## 
## The following object is masked from 'package:car':
## 
##     logit
## 
## 
## The following object is masked from 'package:ggplot2':
## 
##     mean_se
\end{verbatim}

\hypertarget{carga-de-datos-de-parametros-ambientales-de-un-rio}{%
\section{Carga de datos de parametros ambientales de un
Rio}\label{carga-de-datos-de-parametros-ambientales-de-un-rio}}

\hypertarget{dar-formato-a-los-datos}{%
\subsection{Dar formato a los datos}\label{dar-formato-a-los-datos}}

\begin{Shaded}
\begin{Highlighting}[]
\DocumentationTok{\#\#\#Cargar tabla de Datos de parametros para 5 transectos, en los cuales se tomo datos al inicio, medio y final.  }
\CommentTok{\# .csv= formato de valores separados por coma}

\NormalTok{Par\_Transectos }\OtherTok{\textless{}{-}} \FunctionTok{read\_csv}\NormalTok{(}\FunctionTok{here}\NormalTok{(}\StringTok{"Datos"}\NormalTok{,}\StringTok{"Parametros\_Amb.csv"}\NormalTok{))}
\NormalTok{Par\_Transectos2 }\OtherTok{\textless{}{-}} \FunctionTok{read\_csv}\NormalTok{(}\FunctionTok{here}\NormalTok{(}\StringTok{"Datos"}\NormalTok{,}\StringTok{"Parametros\_Amb2.csv"}\NormalTok{))}
\NormalTok{Par\_Transectos3 }\OtherTok{\textless{}{-}} \FunctionTok{read\_csv}\NormalTok{(}\FunctionTok{here}\NormalTok{(}\StringTok{"Datos"}\NormalTok{,}\StringTok{""}\NormalTok{)) }\CommentTok{\# agregar entre comillas el nombre de la tabla}

\CommentTok{\# Explorar los datos, ver los datos de la tabla como hoja de calculo}
\FunctionTok{View}\NormalTok{(Par\_Transectos)}

\CommentTok{\# Ver un resumen de Filas columna y datos contenidos }
\FunctionTok{str}\NormalTok{(Par\_Transectos)}

\DocumentationTok{\#\#\#\#Contraste de Variables entre los 2 muestreos}
\DocumentationTok{\#\#\#Calcular Medidas de Tendencia Central y Dispersion para los parametros fisico{-}quimicos en cada transecto}

\DocumentationTok{\#\#\#Mayo: Crear variable para temperatura ambiental, del agua, conductividad, o2, pH, \% luz, \% humedad, ancho, pendiente para 1° viaje}
\NormalTok{Temp\_Amb1 }\OtherTok{\textless{}{-}}\NormalTok{ Par\_Transectos[,}\DecValTok{2}\NormalTok{]}
\NormalTok{Temp\_Agua1}\OtherTok{\textless{}{-}}\NormalTok{ Par\_Transectos[,}\DecValTok{3}\NormalTok{]}
\NormalTok{Conductividad1}\OtherTok{\textless{}{-}}\NormalTok{ Par\_Transectos[,}\DecValTok{4}\NormalTok{]}
\NormalTok{O21}\OtherTok{\textless{}{-}}\NormalTok{ Par\_Transectos[,}\DecValTok{5}\NormalTok{]}
\NormalTok{pH1}\OtherTok{\textless{}{-}}\NormalTok{ Par\_Transectos[,}\DecValTok{6}\NormalTok{]}
\NormalTok{E\_Luz1}\OtherTok{\textless{}{-}}\NormalTok{ Par\_Transectos[,]}\CommentTok{\# la variable entrada de luz se encuentra en la columna 7}
\NormalTok{Humedad1}\OtherTok{\textless{}{-}}\NormalTok{ Par\_Transectos[,}\DecValTok{8}\NormalTok{]}
\NormalTok{Ancho1}\OtherTok{\textless{}{-}}\NormalTok{ Par\_Transectos[,}\DecValTok{9}\NormalTok{]}
\NormalTok{Profundidad1}\OtherTok{\textless{}{-}}\NormalTok{ Par\_Transectos[,}\DecValTok{10}\NormalTok{]}
\NormalTok{Velocidad1}\OtherTok{\textless{}{-}}\NormalTok{ Par\_Transectos[,}\DecValTok{11}\NormalTok{]}
\NormalTok{Caudal1}\OtherTok{\textless{}{-}}\NormalTok{ Par\_Transectos[,}\DecValTok{12}\NormalTok{]}

\DocumentationTok{\#\#\#Septiembre: Crear variable para temperatura ambiental, del agua, conductividad, o2, pH, \% luz, \% humedad, ancho, pendiente para 2° viaje}
\NormalTok{Temp\_Amb2 }\OtherTok{\textless{}{-}}\NormalTok{ Par\_Transectos2[,}\DecValTok{2}\NormalTok{]}
\NormalTok{Temp\_Agua2}\OtherTok{\textless{}{-}}\NormalTok{ Par\_Transectos2[,}\DecValTok{3}\NormalTok{]}
\NormalTok{Conductividad2}\OtherTok{\textless{}{-}}\NormalTok{ Par\_Transectos2[,}\DecValTok{4}\NormalTok{]}
\NormalTok{O22}\OtherTok{\textless{}{-}}\NormalTok{ Par\_Transectos2[,}\DecValTok{5}\NormalTok{]}
\NormalTok{pH2}\OtherTok{\textless{}{-}}\NormalTok{ Par\_Transectos2[,}\DecValTok{6}\NormalTok{]}
\NormalTok{E\_Luz2}\OtherTok{\textless{}{-}}\NormalTok{ Par\_Transectos2[,}\DecValTok{7}\NormalTok{]}
\NormalTok{Humedad2}\OtherTok{\textless{}{-}}\NormalTok{ Par\_Transectos2[,}\DecValTok{8}\NormalTok{]}
\NormalTok{Ancho2}\OtherTok{\textless{}{-}}\NormalTok{ Par\_Transectos2[,}\DecValTok{9}\NormalTok{]}
\NormalTok{Profundidad2}\OtherTok{\textless{}{-}}\NormalTok{ Par\_Transectos2[,}\DecValTok{10}\NormalTok{]}
\NormalTok{Velocidad2}\OtherTok{\textless{}{-}}\NormalTok{ Par\_Transectos2[,] }\CommentTok{\# Velocidad de corriente se encuentra en columna 11}
\NormalTok{Caudal2}\OtherTok{\textless{}{-}}\NormalTok{ Par\_Transectos2[,}\DecValTok{12}\NormalTok{]}

\DocumentationTok{\#\#\#Octubre: Crear variable para temperatura ambiental, del agua, conductividad, o2, pH, \% luz, \% humedad, ancho, pendiente para 3° viaje}

\DocumentationTok{\#\# Completar Codigo calculando datos octubre o tabla 3}
  
\NormalTok{Temp\_Amb3 }\OtherTok{\textless{}{-}}\NormalTok{ Par\_Transectos3[,}\DecValTok{2}\NormalTok{]}
\NormalTok{Temp\_Agua3}\OtherTok{\textless{}{-}}\NormalTok{ Par\_Transectos3[,}\DecValTok{3}\NormalTok{]}
\NormalTok{Conductividad3}\OtherTok{\textless{}{-}}\NormalTok{ Par\_Transectos3[,}\DecValTok{4}\NormalTok{]}
\NormalTok{O23}\OtherTok{\textless{}{-}}\NormalTok{ Par\_Transectos3[,}\DecValTok{5}\NormalTok{]}
\NormalTok{pH3}\OtherTok{\textless{}{-}}\NormalTok{ Par\_Transectos3[,}\DecValTok{6}\NormalTok{]}
\NormalTok{E\_Luz3}\OtherTok{\textless{}{-}}\NormalTok{ Par\_Transectos3[,}\DecValTok{7}\NormalTok{]}
\NormalTok{Humedad3}\OtherTok{\textless{}{-}}\NormalTok{ Par\_Transectos3[,}\DecValTok{8}\NormalTok{]}
\NormalTok{Ancho3}\OtherTok{\textless{}{-}}\NormalTok{ Par\_Transectos3[,}\DecValTok{9}\NormalTok{]}
\NormalTok{Profundidad3}\OtherTok{\textless{}{-}}\NormalTok{ Par\_Transectos3[,}\DecValTok{10}\NormalTok{]}
\NormalTok{Velocidad3}\OtherTok{\textless{}{-}}\NormalTok{ Par\_Transectos3[,}\DecValTok{11}\NormalTok{] }
\NormalTok{Caudal3}\OtherTok{\textless{}{-}}\NormalTok{ Par\_Transectos3[,}\DecValTok{12}\NormalTok{]}



\DocumentationTok{\#\#\#Calcular medias para parametros por mes}
\DocumentationTok{\#\#Fusionar los Calculos del promedio total y fusionarlos para la tabla de informe}

\DocumentationTok{\#\#Mayo}

\NormalTok{PromTempAmbMay}\OtherTok{\textless{}{-}}\FunctionTok{mean}\NormalTok{(Temp\_Amb1}\SpecialCharTok{$}\StringTok{\textasciigrave{}}\AttributeTok{T° Amb (°C)}\StringTok{\textasciigrave{}}\NormalTok{)}
\NormalTok{PromTemp\_AguaMay}\OtherTok{\textless{}{-}} \FunctionTok{mean}\NormalTok{(Temp\_Agua1}\SpecialCharTok{$}\StringTok{\textasciigrave{}}\AttributeTok{T° Agua (°C)}\StringTok{\textasciigrave{}}\NormalTok{)}
\NormalTok{PromConductividadMay}\OtherTok{\textless{}{-}} \FunctionTok{mean}\NormalTok{(Conductividad1}\SpecialCharTok{$}\StringTok{\textasciigrave{}}\AttributeTok{Conduc (Us/cm)}\StringTok{\textasciigrave{}}\NormalTok{)}
\NormalTok{PromO2May}\OtherTok{\textless{}{-}} \FunctionTok{mean}\NormalTok{(O21}\SpecialCharTok{$}\StringTok{\textasciigrave{}}\AttributeTok{O2 disuel (ppm)}\StringTok{\textasciigrave{}}\NormalTok{)}
\NormalTok{PrompHMay}\OtherTok{\textless{}{-}} \FunctionTok{mean}\NormalTok{(pH1}\SpecialCharTok{$}\NormalTok{pH)}
\NormalTok{PromE\_LuzMay}\OtherTok{\textless{}{-}} \FunctionTok{mean}\NormalTok{(E\_Luz1}\SpecialCharTok{$}\StringTok{\textasciigrave{}}\AttributeTok{\% E. Luz}\StringTok{\textasciigrave{}}\NormalTok{)}
\NormalTok{PromHumedadMay}\OtherTok{\textless{}{-}} \FunctionTok{mean}\NormalTok{(Humedad1}\SpecialCharTok{$}\StringTok{\textasciigrave{}}\AttributeTok{\% Humedad}\StringTok{\textasciigrave{}}\NormalTok{)}
\NormalTok{PromAnchoMay}\OtherTok{\textless{}{-}} \FunctionTok{mean}\NormalTok{(Ancho1}\SpecialCharTok{$}\StringTok{\textasciigrave{}}\AttributeTok{Ancho m}\StringTok{\textasciigrave{}}\NormalTok{)}
\NormalTok{PromProfundidadMay}\OtherTok{\textless{}{-}} \FunctionTok{mean}\NormalTok{(Profundidad1}\SpecialCharTok{$}\StringTok{\textasciigrave{}}\AttributeTok{Prof (cm)}\StringTok{\textasciigrave{}}\NormalTok{)}
\NormalTok{PromVelocidadMay}\OtherTok{\textless{}{-}} \FunctionTok{mean}\NormalTok{(Velocidad1}\SpecialCharTok{$}\StringTok{\textasciigrave{}}\AttributeTok{Vel (m/Seg)}\StringTok{\textasciigrave{}}\NormalTok{)}
\NormalTok{PromCaudalMay}\OtherTok{\textless{}{-}} \FunctionTok{mean}\NormalTok{(Caudal1}\SpecialCharTok{$}\NormalTok{Caudal)}


\DocumentationTok{\#\#Septiembre}

\NormalTok{PromTempAmbSep}\OtherTok{\textless{}{-}}\FunctionTok{mean}\NormalTok{(Temp\_Amb2}\SpecialCharTok{$}\StringTok{\textasciigrave{}}\AttributeTok{T° Amb (°C)}\StringTok{\textasciigrave{}}\NormalTok{)}
\NormalTok{PromTemp\_AguaSep}\OtherTok{\textless{}{-}} \FunctionTok{mean}\NormalTok{(Temp\_Agua2}\SpecialCharTok{$}\StringTok{\textasciigrave{}}\AttributeTok{T° Agua (°C)}\StringTok{\textasciigrave{}}\NormalTok{)}
\NormalTok{PromConductividadSep}\OtherTok{\textless{}{-}} \FunctionTok{mean}\NormalTok{(Conductividad2}\SpecialCharTok{$}\StringTok{\textasciigrave{}}\AttributeTok{Conduc (Us/cm)}\StringTok{\textasciigrave{}}\NormalTok{)}
\NormalTok{PromO2Sep}\OtherTok{\textless{}{-}} \FunctionTok{mean}\NormalTok{(O22}\SpecialCharTok{$}\StringTok{\textasciigrave{}}\AttributeTok{O2 disuel (ppm)}\StringTok{\textasciigrave{}}\NormalTok{)}
\NormalTok{PrompHSep}\OtherTok{\textless{}{-}} \FunctionTok{mean}\NormalTok{(pH2}\SpecialCharTok{$}\NormalTok{pH)}
\NormalTok{PromE\_LuzSep}\OtherTok{\textless{}{-}} \FunctionTok{mean}\NormalTok{(E\_Luz2}\SpecialCharTok{$}\StringTok{\textasciigrave{}}\AttributeTok{\% E. Luz}\StringTok{\textasciigrave{}}\NormalTok{)}
\NormalTok{PromHumedadSep}\OtherTok{\textless{}{-}} \FunctionTok{mean}\NormalTok{(Humedad2}\SpecialCharTok{$}\StringTok{\textasciigrave{}}\AttributeTok{\% Humedad}\StringTok{\textasciigrave{}}\NormalTok{)}
\NormalTok{PromAnchoSep}\OtherTok{\textless{}{-}} \FunctionTok{mean}\NormalTok{(Ancho2}\SpecialCharTok{$}\StringTok{\textasciigrave{}}\AttributeTok{Ancho m}\StringTok{\textasciigrave{}}\NormalTok{)}
\NormalTok{PromProfundidadSep}\OtherTok{\textless{}{-}} \FunctionTok{mean}\NormalTok{(Profundidad2}\SpecialCharTok{$}\StringTok{\textasciigrave{}}\AttributeTok{Prof (cm)}\StringTok{\textasciigrave{}}\NormalTok{)}
\NormalTok{PromVelocidadSep}\OtherTok{\textless{}{-}} \FunctionTok{mean}\NormalTok{(Velocidad2}\SpecialCharTok{$}\StringTok{\textasciigrave{}}\AttributeTok{Vel (m/Seg)}\StringTok{\textasciigrave{}}\NormalTok{)}
\NormalTok{PromCaudalSep}\OtherTok{\textless{}{-}} \FunctionTok{mean}\NormalTok{(Caudal2}\SpecialCharTok{$}\NormalTok{Caudal)}



\DocumentationTok{\#\#Octubre}

\CommentTok{\# completar codigo con calculos para mes de octubre}



\DocumentationTok{\#\#Fusionar los Calculos del promedio total y fusionarlos para la tabla de informe}
\DocumentationTok{\#\#\#Agrupar Los datos del promedio de los parametros por Mes en una tabla}

\NormalTok{PromTempAmb}\OtherTok{\textless{}{-}} \FunctionTok{data.frame}\NormalTok{(PromTempAmbMay, PromTempAmbSep, PromTempAmbOct)}
\NormalTok{PromTempAgu}\OtherTok{\textless{}{-}} \FunctionTok{data.frame}\NormalTok{(PromTemp\_AguaMay, PromTemp\_AguaSep, PromTemp\_AguaOct)}

\NormalTok{PromCondMes}\OtherTok{\textless{}{-}} \FunctionTok{data.frame}\NormalTok{(PromConductividadMay, PromConductividadSep, PromConductividadOct)}

\NormalTok{PromO2Meses}\OtherTok{\textless{}{-}} \FunctionTok{data.frame}\NormalTok{(PromO2May, PromO2Sep, PromO2Oct)}

\NormalTok{PrompHMeses}\OtherTok{\textless{}{-}} \FunctionTok{data.frame}\NormalTok{(PrompHMay, PrompHSep, PrompHOct)}

\NormalTok{PromELuzMes}\OtherTok{\textless{}{-}} \FunctionTok{data.frame}\NormalTok{( PromE\_LuzMay, PromE\_LuzSep, PromE\_LuzOct)}

\NormalTok{PromHumeMes}\OtherTok{\textless{}{-}} \FunctionTok{data.frame}\NormalTok{(PromHumedadMay, PromHumedadSep, PromHumedadOct)}

\NormalTok{PromAnchoMe}\OtherTok{\textless{}{-}} \FunctionTok{data.frame}\NormalTok{(PromAnchoMay, PromAnchoSep, PromAnchoOct)}

\NormalTok{PromProfMes}\OtherTok{\textless{}{-}} \FunctionTok{data.frame}\NormalTok{(PromProfundidadMay, PromProfundidadSep, PromProfundidadOct)}


\NormalTok{PromVelMese}\OtherTok{\textless{}{-}} \FunctionTok{data.frame}\NormalTok{(PromVelocidadMay, PromVelocidadSep, PromVelocidadOct)}

\NormalTok{PromCauMese}\OtherTok{\textless{}{-}} \FunctionTok{data.frame}\NormalTok{(}\CommentTok{\#agregar variable Caudalpara cada mes}


\CommentTok{\# Crea una lista desordenada, por eso hago una lista:}

\FunctionTok{names}\NormalTok{(PromTempAmb)}\OtherTok{\textless{}{-}}\FunctionTok{c}\NormalTok{(}\StringTok{"May"}\NormalTok{, }\StringTok{"Sep"}\NormalTok{,}\StringTok{"Oct"}\NormalTok{)}

\FunctionTok{names}\NormalTok{(PromTempAgu)}\OtherTok{\textless{}{-}} \FunctionTok{c}\NormalTok{(}\StringTok{"May"}\NormalTok{, }\StringTok{"Sep"}\NormalTok{,}\StringTok{"Oct"}\NormalTok{)}


\FunctionTok{names}\NormalTok{(PromCondMes)}\OtherTok{\textless{}{-}} \FunctionTok{c}\NormalTok{(}\StringTok{"May"}\NormalTok{, }\StringTok{"Sep"}\NormalTok{,}\StringTok{"Oct"}\NormalTok{)}
\FunctionTok{names}\NormalTok{(PromO2Meses)}\OtherTok{\textless{}{-}} \FunctionTok{c}\NormalTok{(}\StringTok{"May"}\NormalTok{, }\StringTok{"Sep"}\NormalTok{,}\StringTok{"Oct"}\NormalTok{)}

\FunctionTok{names}\NormalTok{(PrompHMeses)}\OtherTok{\textless{}{-}} \FunctionTok{c}\NormalTok{(}\StringTok{"May"}\NormalTok{, }\StringTok{"Sep"}\NormalTok{,}\StringTok{"Oct"}\NormalTok{)}

\FunctionTok{names}\NormalTok{(PromELuzMes)}\OtherTok{\textless{}{-}} \FunctionTok{c}\NormalTok{(}\StringTok{"May"}\NormalTok{, }\StringTok{"Sep"}\NormalTok{,}\StringTok{"Oct"}\NormalTok{)}

\FunctionTok{names}\NormalTok{(PromHumeMes)}\OtherTok{\textless{}{-}} \FunctionTok{c}\NormalTok{(}\StringTok{"May"}\NormalTok{, }\StringTok{"Sep"}\NormalTok{,}\StringTok{"Oct"}\NormalTok{)}

\FunctionTok{names}\NormalTok{(PromAnchoMe)}\OtherTok{\textless{}{-}} \FunctionTok{c}\NormalTok{(}\StringTok{"May"}\NormalTok{, }\StringTok{"Sep"}\NormalTok{,}\StringTok{"Oct"}\NormalTok{)}

\FunctionTok{names}\NormalTok{(PromProfMes)}\OtherTok{\textless{}{-}} \FunctionTok{c}\NormalTok{(}\StringTok{"May"}\NormalTok{, }\StringTok{"Sep"}\NormalTok{,}\StringTok{"Oct"}\NormalTok{)}


\FunctionTok{names}\NormalTok{(PromVelMese)}\OtherTok{\textless{}{-}} \FunctionTok{c}\NormalTok{(}\StringTok{"May"}\NormalTok{, }\StringTok{"Sep"}\NormalTok{,}\StringTok{"Oct"}\NormalTok{)}


\FunctionTok{names}\NormalTok{(PromCauMese)}\OtherTok{\textless{}{-}} \FunctionTok{c}\NormalTok{(}\StringTok{"May"}\NormalTok{, }\StringTok{"Sep"}\NormalTok{,}\StringTok{"Oct"}\NormalTok{)}


\CommentTok{\# para nombrarlos en vez de [1], uso:}
\FunctionTok{row.names}\NormalTok{(PromTempAmb)}\OtherTok{\textless{}{-}}\FunctionTok{c}\NormalTok{(}\StringTok{"PromTempAmb"}\NormalTok{)}
\FunctionTok{row.names}\NormalTok{(PromTempAgu)}\OtherTok{\textless{}{-}}\FunctionTok{c}\NormalTok{(}\StringTok{"PromTempAgu"}\NormalTok{)}
\FunctionTok{row.names}\NormalTok{(PromCondMes)}\OtherTok{\textless{}{-}}\FunctionTok{c}\NormalTok{(}\StringTok{"PromCondMes"}\NormalTok{)}
\FunctionTok{row.names}\NormalTok{(PromO2Meses)}\OtherTok{\textless{}{-}}\FunctionTok{c}\NormalTok{(}\StringTok{"PromO2Meses"}\NormalTok{)}
\FunctionTok{row.names}\NormalTok{(PrompHMeses)}\OtherTok{\textless{}{-}}\FunctionTok{c}\NormalTok{(}\StringTok{"PrompHMeses"}\NormalTok{)}
\FunctionTok{row.names}\NormalTok{(PromELuzMes)}\OtherTok{\textless{}{-}}\FunctionTok{c}\NormalTok{(}\StringTok{"PromELuzMes"}\NormalTok{)}
\FunctionTok{row.names}\NormalTok{(PromHumeMes)}\OtherTok{\textless{}{-}}\FunctionTok{c}\NormalTok{(}\StringTok{"PromHumeMes"}\NormalTok{)}
\FunctionTok{row.names}\NormalTok{(PromAnchoMe)}\OtherTok{\textless{}{-}}\FunctionTok{c}\NormalTok{(}\StringTok{"PromAnchoMe"}\NormalTok{)}
\FunctionTok{row.names}\NormalTok{(PromProfMes)}\OtherTok{\textless{}{-}}\FunctionTok{c}\NormalTok{(}\StringTok{"PromProfMes"}\NormalTok{)}
\FunctionTok{row.names}\NormalTok{(PromVelMese)}\OtherTok{\textless{}{-}}\FunctionTok{c}\NormalTok{(}\StringTok{"PromVelMese"}\NormalTok{)}
\FunctionTok{row.names}\NormalTok{(PromCauMese)}\OtherTok{\textless{}{-}}\FunctionTok{c}\NormalTok{(}\StringTok{"PromCauMese"}\NormalTok{)}


\DocumentationTok{\#\#\#\#\#\#\#\#\#\#Fusionar las filas de los calculos}

\NormalTok{PromParMesClus}\OtherTok{\textless{}{-}}\FunctionTok{rbind}\NormalTok{(PromTempAmb, PromTempAgu,PromCondMes, PromO2Meses, PrompHMeses, PromELuzMes, PromHumeMes, PromAnchoMe, PromProfMes, PromVelMese, PromCauMese)}

\CommentTok{\#Para obtener resumen de datos en Hoja de calculo.csv}
\FunctionTok{write.csv}\NormalTok{ (PromParMesClus, }\AttributeTok{file=}\StringTok{"PromParMesClus.csv"}\NormalTok{) }\CommentTok{\#buscar en carpeta del proyecto}


\DocumentationTok{\#\#\#Cargar tabla de Datos de parametros de todo la muestra universo setwd("\textasciitilde{}/R/Proyect/Tesis/Parametros/Parametros")}
\NormalTok{Par\_Transectosall}\OtherTok{\textless{}{-}} \FunctionTok{read.csv}\NormalTok{(}\StringTok{"\textasciitilde{}/R/Proyect/Tesis/Parametros/Parametros\_All.csv"}\NormalTok{)}

\DocumentationTok{\#\#\#\#\#\#\#\#\#\#Cluster por mes }

\NormalTok{PromParMesClus}\OtherTok{\textless{}{-}}\FunctionTok{as.data.frame}\NormalTok{(PromParMesClus)}
\DocumentationTok{\#\#LLamar la libreria vegan para hacer cortes en el cladograma}
\FunctionTok{library}\NormalTok{(vegan)}

\DocumentationTok{\#\#\#Realizaremos el analisis de Cluster (clasificaciones gerarquicas), mediante diferentes metodos}
\DocumentationTok{\#\#Metodo de distancia euclidiana por promedios.}
\NormalTok{PromParMesClus}\OtherTok{\textless{}{-}}\FunctionTok{t}\NormalTok{(PromParMesClus) }\CommentTok{\#transponer posicion de filas y columnas en tabla}

\NormalTok{HCMesPar}\OtherTok{\textless{}{-}} \FunctionTok{hclust}\NormalTok{(}\FunctionTok{dist}\NormalTok{(PromParMesClus), }\AttributeTok{method=}\StringTok{"average"}\NormalTok{)}
\CommentTok{\# ver tabla de distancias}
\FunctionTok{summary}\NormalTok{(HCMesPar)}
\FunctionTok{print}\NormalTok{(HCMesPar)}

\CommentTok{\# Graficar analisis}
\FunctionTok{png}\NormalTok{(}\AttributeTok{file=}\StringTok{"HCMesPar.png"}\NormalTok{) }\CommentTok{\# preparando para guardar grafico, omitir esta linea si desea ver el grafico en consola}
\FunctionTok{plot}\NormalTok{(HCMesPar ,}\AttributeTok{hang=}\SpecialCharTok{{-}}\DecValTok{1}\NormalTok{, }\AttributeTok{main=}\StringTok{"Similitud entre Meses de acuerdo a parametros ambientales"}\NormalTok{, }\AttributeTok{labels=}\FunctionTok{c}\NormalTok{(}\StringTok{"Mar"}\NormalTok{, }\StringTok{"Sep"}\NormalTok{, }\StringTok{"Oct"}\NormalTok{))}

\FunctionTok{rect.hclust}\NormalTok{(HCMesPar,}\DecValTok{2}\NormalTok{) }\CommentTok{\# numero de agrupaciones del cluster}
\FunctionTok{dev.off}\NormalTok{() }\CommentTok{\# cerrar lineas y guardar grafico}
\end{Highlighting}
\end{Shaded}

\begin{verbatim}
## Error in parse(text = input): <text>:131:1: unexpected symbol
## 130: 
## 131: names
##      ^
\end{verbatim}

\end{document}
